\subsection{Local Case Studies: From Global Reliance to Individual Forecast Behaviour}
\label{subsec:interp_local_cases}

Global attribution summarizes average model reliance across the held-out
cohort, but it does not show how that reliance is expressed in individual
forecasts. We therefore selected representative local cases to connect the
forecast trajectory with two complementary interpretability views. The left
panel evaluates factual forecast behaviour, the centre panel shows which
historical groups contributed to constructing the streaming state $h_t$, and
the right panel shows which input groups most strongly affected the final
output under leave-one-modality-out ablation. These examples are descriptive
analyses of the trained model and are not estimates of causal physiological
effects.

The left panel in each case shows historical CGM, the forecast median, the
10th--90th percentile interval, the observed future trajectory, the forecast
anchor, and any context annotation such as a sleep window or exercise-detector
interval. The centre panel reports rigorous hidden-state construction
attribution, normalized over the six directly decomposable historical groups:
glucose history, heart rate/physiology, wearable activity/exercise, sleep,
stress, and data quality/missingness. The right panel reports output-level
LOMO reliance, normalized over all eleven input groups visible to the forecast
output. These two percentage scales have different denominators and should not
be compared as numerical ratios. Their rankings and qualitative patterns can
be compared: hidden-state attribution shows what builds the representation,
whereas LOMO shows what changes the output. Neither view proves causal
importance or establishes that a modality improved the forecast in a specific
case.

For each case, we ask three questions at the same forecast anchor: whether the
future glucose trajectory was predicted correctly, which historical modalities
contributed to constructing $h_t$, and which input groups most strongly changed
the final forecast output. The selected anchors are fixed artifacts from the
case-study generation pipeline rather than newly selected examples:
participant 1302, segment 0, anchor 552 for the exercise case; participant
7347, segment 1, anchor 1305 for the stable sleep case; and participant 4162,
segment 0, anchor 651 for the reversal case.

\subsubsection{Exercise-Annotated Forecast-Only Case}
\label{subsec:interp_case_exercise}

Fig.~\ref{fig:interp_case_exercise} shows an exercise-annotated factual
forecast. The Study 2 exercise detector was external to the base forecasting
model. It was used to identify and annotate the walking episode, but the
detector-active state was not a dedicated input to the displayed forecast.
Historical steps, heart rate, and related wearable measurements were available
to the model through the standard wearable and physiology channels. This case
therefore evaluates the forecast-only pathway, not the exercise counterfactual
branch. The hatched interval marks an independently detected 70-minute walking
bout with moderate strain and is displayed for interpretation; it is not a
forecast input under the checkpoint inspection summarized in
\texttt{external\_candidate\_check\_v2.csv}.

\begin{figure}[H]
    \centering
    \includegraphics[width=\linewidth]{generated/fig_case_exercise_v2.png}
    \caption{%
      \textbf{Exercise-annotated forecast-only case.}
      Participant 1302, segment 0, anchor 552.
      The episode was independently identified by the Study 2 detector, but the
      detector state, cadence and event phase were not supplied as dedicated
      inputs to the displayed forecast. Left: historical CGM, forecast median,
      10th to 90th percentile interval, observed future and detector-annotated
      exercise interval. Centre: rigorous hidden-state construction attribution
      normalized over the six directly decomposable historical groups. Right:
      output-level LOMO reliance normalized over all eleven input groups.
      Strong activity representation in $h_t$ does not imply that an explicit
      exercise-response function was available to the forecast decoder.%
    }
    \label{fig:interp_case_exercise}
\end{figure}

The forecast in Fig.~\ref{fig:interp_case_exercise} starts from 90~mg/dL and
captures the overall post-anchor recovery toward higher glucose, although the
observed trajectory remains below the forecast median through most of the
hour. The horizon-average MAE is 2.95~mg/dL, the 60-minute absolute error is
1.78~mg/dL, and the maximum absolute error is 4.74~mg/dL at 40 minutes. All 12
observations remain inside the 10th--90th percentile interval, whose width is
30.52~mg/dL at 60 minutes. This is therefore an accurate, interval-covered
forecast, albeit with a relatively broad uncertainty band.

The state-construction view shows that wearable activity/exercise is the
largest contributor to $h_t$ in this case, accounting for 27.6\% of the
six-group hidden-state attribution. Data quality/missingness contributes
22.5\%, glucose history 18.3\%, and heart rate/physiology 14.1\%. This confirms
that raw wearable activity information reached the streaming representation.
It does not prove that the model constructed an explicit semantic state for
exercise onset, duration, cadence, recovery phase, or time since exercise. In
the output-level view, the final forecast is dominated by glucose history
(50.2\%), followed by data quality/missingness (14.4\%), sleep (7.7\%), and
time/calendar (7.6\%). Wearable activity/exercise contributes only 2.4\% of
the eleven-group LOMO attribution, while heart rate/physiology contributes
5.3\%.

The episode was independently identified by the Study 2 detector as a
70-minute walking bout with moderate strain. Activity-related historical inputs
were strongly represented in $h_t$, and the factual forecast was accurate, but
the displayed trajectory was produced in forecast-only mode: detector state,
exercise onset, cadence class, recovery phase, and time since exercise were not
provided as dedicated decoder inputs. The successful trajectory should
therefore be interpreted as performance of the general multimodal forecaster
for this selected episode, not as evidence that the model learned an explicit
exercise-response function.

Because the external detector uses steps and heart rate, elevated activity
attribution in a detector-selected exercise episode is partly expected by
construction. The informative result is the separation between strong activity
representation (27.6\% of $h_t$ construction) and weak direct output reliance
(2.4\% under LOMO), despite an accurate forecast. This local case evaluates
factual forecasting and representation. The separate exercise-counterfactual
analysis evaluates how the decoder output changes when an activity scenario is
explicitly asserted.

\subsubsection{Stable Overnight Sleep Forecast Case}
\label{subsec:interp_case_sleep}

Fig.~\ref{fig:interp_case_sleep} provides a second well-calibrated local case
before the reversal failure case. Glucose is stable at the anchor: the current glucose is
108~mg/dL and the 30-minute pre-anchor slope is 0.033~mg/dL/min. The model
predicts a nearly flat trajectory, the observed future declines only slightly,
and the observed trajectory remains inside the prediction interval. The
60-minute absolute error is 3.28~mg/dL, the horizon-average MAE is
1.95~mg/dL, interval coverage is 100.0\%, and the 60-minute interval width is
8.51~mg/dL. This is a small-error, well-calibrated forecast.

\begin{figure}[H]
    \centering
    \includegraphics[width=\linewidth]{generated/fig_case_stable_sleep_v2.pdf}
    \caption{%
      \textbf{Stable overnight sleep forecast case.}
      Left: sleep window, historical CGM, forecast median, prediction interval
      and observed future. Centre: hidden-state construction attribution over
      the six historical groups. Right: output-level LOMO reliance over all
      eleven groups. The balanced multimodal state is translated into a stable
      forecast whose observed future remains inside the prediction interval.%
    }
    \label{fig:interp_case_sleep}
\end{figure}

The hidden-state attribution in Fig.~\ref{fig:interp_case_sleep} is multimodal
rather than CGM-only. Data quality/missingness is the largest
state-construction group at 27.3\%, followed by heart rate/physiology at
22.4\%, glucose history at 17.0\%, and sleep at 16.4\%. The output-level LOMO
view remains led by glucose history at 29.2\%, but site/cohort (14.1\%),
heart rate/physiology (11.8\%), data quality/missingness (10.1\%), and sleep
(9.9\%) also influence the output. Compared with the exercise case, sleep and
physiology are more visibly translated into the final forecast.

The stable overnight case illustrates a context in which multimodal state
information is translated into an accurate and calibrated forecast. Data
quality, heart-rate physiology, glucose history and sleep all contribute to
construction of the hidden state, while the final output remains primarily
glucose-driven. The observed future remains within the prediction interval and
differs only slightly from the forecast median. In this relatively stable
context, the model successfully combines recent glucose with sleep and
physiological information.

The large contribution of data-quality variables should not be interpreted as
a physiological mechanism. It may reflect legitimate reliability conditioning,
but it may also encode sleep state, wearable non-wear, participant identity,
site workflow or device-specific missingness.

\subsubsection{Post-Decline Reversal Failure Case}
\label{subsec:interp_case_reversal}

The reversal example in Fig.~\ref{fig:interp_case_reversal} uses the
\texttt{fast\_down} case rendered as \texttt{fig\_case\_reversal\_v2}. The
case-quality audit confirms that the future CGM values are observed, no
invalid CGM gap crosses the forecast horizon, no participant or segment
boundary is crossed, timestamps are correctly aligned, no inverse-scaling error
exists, values remain inside the accepted CGM sensor range, and the trajectory
is not generated by a preprocessing artifact. The corresponding audit row has
\texttt{all\_checks\_pass=True}.

\begin{figure}[t]
    \centering
    \includegraphics[width=\linewidth]{generated/fig_case_reversal_v2.pdf}
    \caption{%
      \textbf{Post-decline reversal failure.}
      The model extrapolates the recent downward glucose tendency, whereas
      observed glucose reverses sharply and exceeds the upper prediction
      interval. Glucose history dominates both state construction and output
      reliance. The example illustrates directional, magnitude and uncertainty
      failure under an abrupt trajectory change.%
    }
    \label{fig:interp_case_reversal}
\end{figure}

Glucose declines strongly before the forecast anchor, with a 30-minute
pre-anchor slope of -3.300~mg/dL/min and current glucose of 169~mg/dL. The
model predicts a small further decline followed by a relatively stable level,
whereas the observed future reverses sharply and rises above 300~mg/dL. The
observed trajectory exceeds the upper prediction interval rapidly: interval
coverage is 0.0\%, the first interval exit occurs at 5 minutes, the maximum
distance outside the interval is 84.76~mg/dL, and the 60-minute absolute error
is 136.33~mg/dL.

This failure has three components. It is a directional failure because the
forecast continues the recent decline rather than predicting the reversal. It
is a magnitude failure because the later increase is strongly underestimated.
It is an uncertainty failure because the observed trajectory lies outside the
prediction interval throughout the horizon.

The centre panel shows that glucose history is the largest contributor to
$h_t$, accounting for 40.7\% of the six-group hidden-state attribution.
Activity is also represented at 19.8\%, and heart rate/physiology at 17.8\%.
The recent declining glucose trajectory is therefore the dominant
state-construction signal. In the right panel, glucose history also dominates
final-output reliance at 33.1\%. Static clinical information contributes
14.2\%, wearable activity/exercise 13.7\%, site/cohort 11.6\%, and
time/calendar 11.4\%. Heart rate/physiology is represented in $h_t$ but has
little direct output influence in this case, contributing 1.9\% of the
eleven-group LOMO attribution.

The reversal case exposes a major limitation of the forecast-only model. The
sharply declining pre-anchor glucose trajectory dominates both hidden state
construction and final-output reliance, and the decoder continues the recent
downward tendency. The observed future instead reverses rapidly and rises well
above the upper prediction interval. Activity and heart-rate information are
present in the state, but they are not translated into a prediction of the
reversal. This example therefore represents a directional, magnitude and
uncertainty-calibration failure.

The reversal may depend on information that cannot be identified from the
recent observed trajectory, such as an unrecorded meal, treatment response or
another physiological disturbance. No specific cause should be assigned without
a corresponding logged event.

\subsubsection{What the Local Cases Reveal}
\label{subsec:interp_case_synthesis}

The three cases separate representation from response modelling. In the
exercise episode, wearable activity strongly contributes to the hidden state,
but the decoder does not fully predict the later glucose decline. In the
stable overnight case, multimodal state information is translated into an
accurate and calibrated forecast. In the post-decline reversal, recent glucose
history dominates both the state and output, and the decoder continues the
recent trend while missing an abrupt change in direction. These cases suggest
that the encoder captures meaningful multimodal context, while the main
remaining limitation is the decoder's ability to convert that context into
event-specific future glucose trajectories.

The exercise case motivates explicit event-aware modelling rather than relying
only on raw wearable history. A future extension could condition a residual
head or scenario decoder on exercise-active state, time since onset, time since
end, duration, cadence, strain and participant phenotype. The current case
study motivates this extension but does not evaluate it.

The local examples are selected illustrations of model behaviour. They support
interpretation of the cohort-level analyses but do not establish the frequency
of each failure mode. Population-level conclusions must rely on systematic
stratified evaluation across all eligible anchors.
